\section{Documentary dependency audit of the proposed global-to-local calculation}
\label{sec:source-dependency-audit}

This section records a fresh complete reading of Masullo~(2026),
\S\S12.3--12.5, printed pp.~514--520.  It compares the exact
source assertions with results already proved in this report.
It introduces no new mathematical theorem or arithmetic conclusion.

\paragraph{The absent comparison.}
Section~12.3, p.~514, has its heading and the footnote
``Include section'', but no argument.  On p.~516, Step~III
explicitly requires an identification of the local truncated
kernel with the kernel denoted by $k^{\mathrm{PW}}$.
On p.~517, the proof computes the kernel of a local physical
Fourier-cutoff composition and assigns it that notation.
This computation does not identify that kernel with the one
defined using the weighted global orthogonal PW projection.

MA-PI-001--003 refute a precise canonical theta intertwining
between the two constructions on actual source vectors.
MA-PI-002 uses the literal printed Hilbert realization,
including the norm-half twist specified on p.~515;
MA-PI-003 also retains both polar conditions in the principal
sector.  MA-PI-004 supplies the non-induced cubic character
over $\mathbb Q(\zeta_3)$, retaining the different prime over
$3$, the conductor prime over $269$, both physical Fourier
shells and genuine flat source data.
The draft does not print the exact canonical equation
tested by these results.  The documentary conclusion is therefore
that the required identification is absent and this particular
candidate fails.  No claim about every different factorization,
or every scalar trace equality, follows.

\paragraph{Hypotheses and cutoffs on p.~514.}
Theorem~12.4.1.1 retains class-group triviality after localization,
all archimedean places, all places ramified over $\mathbb Q$,
and the indicated outside-$S$ unit invariance of the test.
Field ramification cannot be replaced merely by character
ramification.  Footnote~2 flags unfinished work on a one-sided
support convention, an independent lower cutoff, and the resulting
limit.  The actual repaired local quotient operators are supplied
in MA-LA-014--015, MA-CQ-003--004 and MA-GA-001 at their stated
scopes.  The frequency operator in Step~II, p.~516, is explicitly
a product-norm ball.  The annular-frequency alternative of
MA-LA-016 is a different operator.

\paragraph{The source and trace factorization on p.~515.}
The proof fixes the norm-half twisted Sobolev realization,
asserts nuclearity on a quotient PW level, and invokes a
factorization through a local theta image before cyclic trace
transfer.  Its topological-copy assertion cites
Proposition~5.3.4.4 and is accompanied by a footnote asking for
that proposition to be updated.  The cited transform/image
criterion does not itself specify the continuous inverse,
nuclear lift, or commuting operator diagram needed here.
Range inclusion alone does not justify cyclic transfer.
The actual local sections and returns in MA-LA-011, MA-CQ,
MA-GA and MA-MO-017/028 have their proved domains and metrics;
they do not by themselves identify the global orthogonal
projection with the physical Fourier cutoff.

\paragraph{Literal assertions decided by the earlier source audit.}
MA-LA-008 identifies the output norm-half factor present in the
p.~515 operator but absent from the raw $K$ in (12.4.1.7),
p.~516.  The exact repair is $A=UK$; a change of coordinates
also requires consistent conjugation of the other operators.
MA-LA-009 refutes the raw tensor containment following that
display: equivariance does not make the raw output constant on
input $S$-unit coinvariant classes.  An output quotient or a
specified section is required.
MA-LA-010 refutes the asserted additive-Schwartz domain/codomain
of mixed-place product-norm cutoffs, even after an annular
position repair.  They remain bounded additive $L^2$ multipliers;
that weaker fact does not establish the printed source mapping.

\paragraph{Local kernels and nuclearity on pp.~516--518.}
Equations~(12.4.1.8)--(12.4.1.10) compute the local physical
composition.  Its kernel and the ensuing local orbital-volume
calculation still require the actual quotient/source domains.
MA-LA-012 rules out interpreting a nonzero unit-equivariant
operator of this form as compact on the full additive $L^2$
cover.  This is an attribution boundary: the p.~515 nuclearity
claim concerns a quotient PW level, and p.~517 uses a fundamental
domain and unit sum.  We do not attribute an explicit full-cover
trace-class theorem to the draft.  Actual quotient trace-class
operators exist and have been constructed above.
The p.~518 logarithmic-volume calculation supplies neither
the missing global map nor the lower-cutoff/rate analysis
flagged in footnote~2.

\paragraph{Rates and the final reuse on pp.~519--520.}
Page~519 requires both smoothing-width products to vanish,
but does not impose the additional slow-band relation between
the large radii.  MA-LA-015 proves profile comparison for the
actual repaired local operators; its paired global comparison
retains the stated slow-band hypothesis.
Theorem~12.5.1.1 assumes vanishing on the entire norm-one
subgroup.  The p.~520 choice of $S$ retains field and character
ramification, structural/class-group conditions, test invariance,
and all relevant small primes.

The p.~520 argument first decomposes the global kernel into
its $S$-part and outside-$S$ parts, then removes the latter
using the cited shell theorem.  It next invokes
Theorem~12.4.1.1 to identify the global $S$-kernel trace with
the local distributions minus the two polar terms, and finally
uses the spectral construction to obtain the relative-strip
zero sum.  The third step is the unsupplied comparison.
The local calculation alone gives neither equality of these
kernels, compatibility of their Hilbert metrics, control of
the limiting remainder, nor the polar subtraction.
The missing \S12.3 is the intended location of such a bridge.

\paragraph{Scope of the resulting conclusion.}
The source contains literal repairable errors, assertions
refuted on their stated domains, a local calculation for which
actual quotient repairs are now proved, and a separate missing
global comparison.  These categories must not be conflated.
MA-FT-001/002 and the Gram/resolvent results retain their exact
operator limits and arithmetic deficiency condition.
The canonical-map counterexamples do not eliminate that
condition or rule out all other comparison maps.
This documentary pass uses the previously recorded whole-proof
reviews of those results; it does not claim a fresh complete
proof review of every cited module.  Its row-level evidence is
in \texttt{source\_dependency\_audit.tsv}.
